\section{Related work}\label{sec:related}

\paragraph{Unpivoted LU for nonsingular matrices}
The classical existence theory assumes invertibility.  If $A$ is nonsingular,
then $A = LU$ exists if and only if all leading principal submatrices are
nonsingular, in which case $L$ and $U$ are invertible and the factorization
with unit lower triangular $L$ is unique; see
Higham~\cite{higham_gaussian_2011} for a survey and
\cite{golub_matrix_2013} for textbook treatments.  Higham notes explicitly
that when a proper leading principal submatrix is singular a factorization
"may exist, but if so it is not unique," without characterizing when.  Lau
and Markham~\cite{lau_markham_1979} and Johnson, Olesky, and van den
Driessche~\cite{johnson_inherited_1989} showed that a possibly singular $A$
has a \emph{unique} unit LU factorization precisely when the $n-1$
\emph{proper} leading principal minors are nonzero; the analysis
of~\cite{johnson_inherited_1989} (inheritance of entries, fill-in) assumes
this condition throughout and does not address existence when it fails.

\paragraph{Unpivoted LU for arbitrary matrices}
The existence question for arbitrary, possibly singular $A$ was settled by
Okunev and Johnson.  Their result, circulated as a 1997 manuscript and posted
as~\cite{okunev_johnson_2005}, recently appeared
as~\cite{johnson_okunev_2025}: $A \in \F^{n \by n}$ admits $A = LU$, with no
invertibility requirement on the factors, if and only if
\begin{equation}\label{eq:oj}
  \rk(A[1:k,1:k]) + k \ \ge\ \rk(A[1:k,:]) + \rk(A[:,1:k]),
  \qquad k = 1, \ldots, n .
\end{equation}
Their sufficiency proof is an induction on "factor complements"
(generalized Schur complements), choosing the first column of $L$ and the
first row of $U$ case by case; the published version also develops
"near-LU" factorizations with almost-triangular factors when \cref{eq:oj}
fails, and re-derives $PLU$ and three-factor decompositions.
Darve~\cite{darve_lu_2026} restates \cref{eq:oj} in the nullity form used in
this paper, gives a different constructive proof based on \emph{restricted
pivoting} (only rows or columns that are zero in the active Schur complement
are moved, so that undoing the internal permutations preserves
triangularity), and develops rank-revealing structure, tight sparsity bounds
for the factors, and criteria for one-sided unit-triangular factors.  The
present paper gives a third, normal-form proof of the same criterion, whose
main advantage is that it extends verbatim to the symmetric $LDL^T$ problem.

Dopico, Johnson, and Molera~\cite{dopico_multiple_2006} study the related
question in which $L$ is required to be unit lower triangular (hence
invertible): existence is characterized through a canonical form of $A$
under left multiplication by unit lower triangular matrices (the
rank-revealing minimal canonical form), and, importantly, all LU
factorizations of a singular matrix are parametrized, quantifying their
non-uniqueness.  Our normal form differs in that the reduction acts on both
sides (or by congruence), and the middle factor --- not an upper triangular
matrix --- carries the combinatorial obstruction.
Nagarajan, Devasahayam, and Soundararajan~\cite{nagarajan_products_1999}
showed that \emph{every} square matrix over a field is a product of three
triangular matrices $LUM$; the two-factor existence question is exactly the
question of when the third factor can be dispensed with.
Jeffrey~\cite{jeffrey_lu_2010} standardizes full-rank LU factorings of
rectangular and rank-deficient matrices in the form $P_r L D^{-1} U P_c$;
since row and column permutations are built in, existence is unconditional
and the unpivoted question does not arise.  For the structured class of
M-matrices, McDonald and Schneider~\cite{mcdonald_schneider_1998} gave
graph-theoretic necessary and sufficient conditions for a singular M-matrix
to admit an unpivoted (block) LU factorization into possibly singular
M-matrix factors.  Rank-revealing LU with pivoting chosen for numerical
size, rather than algebraic existence, is developed by Miranian and
Gu~\cite{miranian_gu_2003}, and the conditioning of the triangular factors
of pivoted factorizations is analyzed by
Stewart~\cite{stewart_triangular_1997}.  General surveys of matrix
factorizations include~\cite{lu_numerical_2021}.

\paragraph{Symmetric factorizations}
For symmetric positive semidefinite matrices over $\mathbb{R}$, an unpivoted
Cholesky-type factorization always exists even in the singular case: $A =
LL^T$ with $L$ lower triangular, possibly with zero diagonal entries; see
Higham~\cite{higham_cholesky_1990} for a detailed treatment.  Our symmetric
criterion is consistent with this: positive semidefiniteness forces
$\nl(A[1:k,1:k]) = \nl(A[:,1:k])$ for all $k$, so the criterion of
\Cref{thm:ldlt} holds automatically (\Cref{cor:psd}).  For symmetric
\emph{indefinite} matrices the classical approach abandons the diagonal $D$:
with symmetric pivoting, $PAP^T = LDL^T$ always exists when $D$ is allowed
to be block diagonal with $1 \by 1$ and $2 \by 2$ blocks, as introduced by
Bunch and Parlett~\cite{bunch_parlett_1971} and made efficient by Bunch and
Kaufman~\cite{bunch_kaufman_1977}; a rank-profile-revealing variant over
arbitrary fields, still with symmetric pivoting and with $2 \by 2$
antidiagonal blocks in $D$, is given by Dumas and
Pernet~\cite{dumas_pernet_2018}.  Closest in spirit to our symmetric
reduction is the elegant paper of Van Dooren and
Strang~\cite{vandooren_strang_2014}, which factors a real symmetric matrix
as $A = L P L^T$ \emph{without pivoting}, where $P$ is a signed-matching
middle factor built from diagonal entries $0, \pm 1$ and symmetric exchange
blocks; our \Cref{lem:reduction-sym} can be read as a field-general,
weighted version of their reduction (\Cref{rem:signed}).  Neither
\cite{vandooren_strang_2014} nor \cite{dumas_pernet_2018} addresses when
the middle factor can be made \emph{truly diagonal}, which is exactly the
step from $A = MPM^T$ to $A = LDL^T$ taken in
\Cref{sec:combinatorics,sec:explicit}.  The $2 \by 2$ blocks of
Bunch--Parlett reappear in our theory in a different role: they are the
irreducible blocks of the symmetric matching normal form, i.e., the
obstruction that must be threaded through earlier zero indices rather than a
device enabled by pivoting.  We are not aware of a previously published
necessary and sufficient condition for the unpivoted diagonal-$D$ problem
with possibly singular $L$.

\paragraph{Triangular canonical forms and the Bruhat decomposition}
Our reduction lemma states that every $A$ equals $M \Pi N$ with $M, N$
invertible triangular and $\Pi$ a matching matrix, and that the matching is
an invariant of $A$.  For invertible $A$ this is a form of the Bruhat
decomposition $A = LPU$, whose permutation $P$ is uniquely determined by
$A$; see Elsner~\cite{elsner_bruhat_1979} and
Strang~\cite{strang_pnas_2010,strang_lpu_2012} for elimination-based
accounts, and~\cite{dopico_multiple_2006} for the related one-sided
canonical form.  For singular matrices the permutation becomes a partial
permutation: this is the generalized Bruhat (or LEU)
decomposition~\cite{malaschonok_2010}, and the invariant matching is the
\emph{rank profile matrix} of Dumas, Pernet, and
Sultan~\cite{dumas_pernet_sultan_2017}, with which part~(1) of our
\Cref{prop:canonical} coincides.  On the symmetric side,
Szechtman~\cite{szechtman_2007} classified symmetric matrices under
congruence by triangular groups: for $\operatorname{char} \F \neq 2$ the
canonical forms under unit-triangular congruence are weighted symmetric
sub-permutation matrices over any field, while over fields of
characteristic $2$ in which every element is a square the canonical forms
must be relaxed to include the $2 \by 2$ blocks with a nonzero corner that
appear in our \Cref{def:sym-matching}; see also Bagno and
Cherniavsky~\cite{bagno_cherniavsky_2012} for the rank-control invariants
of congruence $B$-orbits over $\mathbb{C}$.  Our
\Cref{lem:reduction-sym,prop:canonical} thus recover known canonical forms;
we include short elimination-style constructive proofs both to keep the
paper self-contained and because the constructions are exactly what the
factorization algorithm of \Cref{sec:validation} executes.

\paragraph{Positioning}
In summary: the unsymmetric criterion in \Cref{thm:lu} is due to Okunev and
Johnson~\cite{okunev_johnson_2005,johnson_okunev_2025}, with alternative
proofs and structural refinements in~\cite{darve_lu_2026}; the triangular
canonical forms we use connect to the generalized Bruhat decomposition
\cite{elsner_bruhat_1979,malaschonok_2010,dumas_pernet_sultan_2017} and, in
the symmetric case, to Van Dooren--Strang~\cite{vandooren_strang_2014} and
Szechtman~\cite{szechtman_2007}.  The contributions of the present paper
are the unified normal-form route to \emph{existence} theorems, the
symmetric $LDL^T$ criterion (\Cref{thm:ldlt}) together with the equivalence
of $LU$ and $LDL^T$ existence for symmetric matrices
(\Cref{cor:sym-lu-ldlt}), the prefix-condition combinatorics and threading
constructions of \Cref{sec:combinatorics,sec:explicit}, and the
exhaustively validated constructive algorithm.
